\chapter{Hito 4: Diseño del Modelo Dimensional}

\section{Objetivo}

El objetivo de este hito es diseñar el modelo dimensional del Data Warehouse macrofinanciero.

Al finalizar este hito el estudiante será capaz de:

\begin{itemize}
\item Identificar el grano de cada tabla de hechos.
\item Diferenciar dimensiones y hechos.
\item Diseñar un esquema estrella.
\item Definir claves de negocio y claves subrogadas.
\item Documentar el modelo antes de implementarlo en PostgreSQL.
\end{itemize}

\section{Resultado esperado}

Al finalizar este hito tendremos definido el modelo lógico del Data Warehouse.

El modelo estará compuesto por:

\begin{itemize}
\item \texttt{dim\_tiempo}
\item \texttt{dim\_moneda}
\item \texttt{dim\_entidad\_financiera}
\item \texttt{fact\_tipo\_cambio}
\item \texttt{fact\_balance\_mensual}
\end{itemize}

\section{Paso 1: Entender qué es un modelo dimensional}

Un modelo dimensional organiza los datos para facilitar el análisis.

Está compuesto por dos tipos principales de tablas:

\begin{itemize}
\item \textbf{Dimensiones}: describen el contexto del análisis.
\item \textbf{Hechos}: contienen medidas numéricas que queremos analizar.
\end{itemize}

Ejemplo:

\begin{itemize}
\item La fecha es una dimensión.
\item La entidad financiera es una dimensión.
\item La moneda es una dimensión.
\item La tasa de cambio es una medida.
\item Los activos bancarios son una medida.
\end{itemize}

\section{Paso 2: Identificar los procesos de negocio}

En este proyecto tendremos dos procesos de negocio principales:

\begin{enumerate}
\item Seguimiento del mercado cambiario.
\item Seguimiento de balances bancarios mensuales.
\end{enumerate}

Cada proceso generará una tabla de hechos.

\section{Paso 3: Definir el grano}

El grano indica qué representa una fila de una tabla de hechos.

Esta decisión es fundamental.

\subsection{Grano de fact\_tipo\_cambio}

\begin{quote}
Una fila por fecha y moneda.
\end{quote}

Ejemplo:

\begin{verbatim}
2024-01-01 | USD
2024-01-01 | EUR
2024-01-02 | USD
\end{verbatim}

\subsection{Grano de fact\_balance\_mensual}

\begin{quote}
Una fila por mes y entidad financiera.
\end{quote}

Ejemplo:

\begin{verbatim}
2024-01-31 | B001
2024-01-31 | B002
2024-02-29 | B001
\end{verbatim}

\section{Paso 4: Definir dimensiones}

\subsection{Dimensión tiempo}

La dimensión tiempo será compartida por ambos hechos.

Esto la convierte en una dimensión conformada.

\begin{verbatim}
dim\`_tiempo

sk\_tiempo
fecha
anio
mes
anio\_mes
nombre\_mes
trimestre
es\_fin\_de_mes
\end{verbatim}

\subsection{Dimensión moneda}

Esta dimensión describe las monedas utilizadas en el mercado cambiario.

\begin{verbatim}
dim\_moneda

sk\_moneda
moneda\_cod
moneda\_nombre
\end{verbatim}

\subsection{Dimensión entidad financiera}

Esta dimensión describe las entidades financieras.

\begin{verbatim}
dim\_entidad\_financiera

sk\_entidad
entidad\_cod
entidad\_nombre
tipo\_entidad
vigente\_desde
vigente\_hasta
es\_actual
\end{verbatim}

La dimensión entidad financiera se diseñará para soportar cambios históricos mediante SCD Tipo 2.

\section{Paso 5: Definir tablas de hechos}

\subsection{Hecho tipo de cambio}

\begin{verbatim}
fact\_tipo\_cambio

sk\_tiempo
sk\_moneda
tasa_compra
tasa_venta
tasa_promedio
spread
\end{verbatim}

Medidas:

\begin{itemize}
\item \texttt{tasa\_compra}
\item \texttt{tasa\_venta}
\item \texttt{tasa\_promedio}
\item \texttt{spread}
\end{itemize}

\subsection{Hecho balance mensual}

\begin{verbatim}
fact\_balance\_mensual

sk\_tiempo
sk\_entidad
activos\_totales
cartera\_creditos
depositos
patrimonio
resultado\_neto
depositos\_me
cartera\_me
\end{verbatim}

Medidas:

\begin{itemize}
\item \texttt{activos\_totales}
\item \texttt{cartera\_creditos}
\item \texttt{depositos}
\item \texttt{patrimonio}
\item \texttt{resultado\_neto}
\item \texttt{depositos\_me}
\item \texttt{cartera\_me}
\end{itemize}

\section{Paso 6: Definir claves}

\subsection{Clave de negocio}

La clave de negocio proviene de la fuente original.

Ejemplos:

\begin{itemize}
\item \texttt{moneda\_cod}: USD, EUR.
\item \texttt{entidad\_cod}: B001, B002.
\item \texttt{fecha}\: 2024-01-31.
\end{itemize}

\subsection{Clave subrogada}

La clave subrogada es creada por el Data Warehouse.

Ejemplos:

\begin{itemize}
\item \texttt{sk\_tiempo}
\item \texttt{sk\_moneda}
\item \texttt{sk\_entidad}
\end{itemize}

Las tablas de hechos no deben depender directamente de las claves de negocio.

Deben conectarse mediante claves subrogadas.

\section{Paso 7: Diseñar el esquema estrella}

El esquema estrella del proyecto será:

\begin{verbatim}
dim_tiempo
|
|
dim_moneda ---- fact_tipo_cambio

    dim_tiempo
       |
       |

dim_entidad_financiera ---- fact_balance_mensual
\end{verbatim}

La dimensión \texttt{dim\_tiempo} conecta ambos hechos y permite análisis integrados.

\section{Paso 8: Crear documento del modelo}

Crear el archivo:

\begin{verbatim}
docs/modelo\_dimensional.md
\end{verbatim}

Contenido sugerido:

\begin{verbatim}

# Modelo Dimensional

## Procesos de negocio

1. Mercado cambiario
2. Balances bancarios mensuales

## Tablas de hechos

### fact_tipo_cambio

Grano:
Una fila por fecha y moneda.

Medidas:

* tasa_compra
* tasa_venta
* tasa_promedio
* spread

### fact_balance_mensual

Grano:
Una fila por mes y entidad financiera.

Medidas:

* activos_totales
* cartera_creditos
* depositos
* patrimonio
* resultado_neto
* depositos_me
* cartera_me

## Dimensiones

* dim_tiempo
* dim_moneda
* dim_entidad_financiera
  \end{verbatim}

\section{Paso 9: Crear tabla de definición del modelo}

Agregar al documento una tabla como la siguiente:

\begin{center}
\begin{tabular}{lll}
\toprule
Tabla & Tipo & Grano \\
\midrule
dim\_tiempo & Dimensión & Una fila por día \\
dim\_moneda & Dimensión & Una fila por moneda \\
dim\_entidad\_financiera & Dimensión & Una fila por versión de entidad \\
fact\_tipo\_cambio & Hecho & Una fila por fecha y moneda \\
fact\_balance\_mensual & Hecho & Una fila por mes y entidad \\
\bottomrule
\end{tabular}
\end{center}

\section{Paso 10: Definir indicadores derivados}

A partir de este modelo construiremos indicadores como:

\begin{itemize}
\item Tipo de cambio promedio.
\item Spread cambiario.
\item ROA.
\item ROE.
\item Dolarización de depósitos.
\item Dolarización de cartera.
\item Concentración bancaria.
\end{itemize}

Fórmulas principales:

\begin{verbatim}
tasa_promedio = (tasa_compra + tasa_venta) / 2

spread = tasa_venta - tasa_compra

ROA = resultado_neto / activos_totales

ROE = resultado_neto / patrimonio

dolarizacion_depositos = depositos_me / depositos

dolarizacion_cartera = cartera_me / cartera_creditos
\end{verbatim}

\section{Paso 11: Validar el diseño}

Antes de crear las tablas en PostgreSQL, responder:

\begin{enumerate}
\item ¿Cada tabla de hechos tiene un grano claro?
\item ¿Cada medida pertenece al grano definido?
\item ¿Las dimensiones describen correctamente los hechos?
\item ¿La dimensión tiempo puede usarse en ambos hechos?
\item ¿La entidad financiera requiere historia?
\end{enumerate}

\section{Checklist de validación}

Antes de continuar verificar:

\begin{itemize}
\item[$\Box$] Procesos de negocio identificados.
\item[$\Box$] Grano de cada hecho definido.
\item[$\Box$] Dimensiones definidas.
\item[$\Box$] Hechos definidos.
\item[$\Box$] Claves de negocio identificadas.
\item[$\Box$] Claves subrogadas definidas.
\item[$\Box$] Documento \texttt{modelo\_dimensional.md} creado.
\end{itemize}

\section{Entregable}

Al finalizar este hito el estudiante deberá disponer de:

\begin{itemize}
\item Diseño lógico del Data Warehouse.
\item Documento del modelo dimensional.
\item Definición de grano.
\item Lista de dimensiones.
\item Lista de hechos.
\item Lista de indicadores derivados.
\end{itemize}

\section{Próximo hito}

En el Hito 5 implementaremos la primera capa física del Data Warehouse:

\begin{itemize}
\item Creación de schemas.
\item Creación de tablas staging.
\item Preparación de la capa bronze.
\item Primer script SQL del modelo.
\end{itemize}
